% !TEX root = ../main.tex
\section{Introduction}

Stochastic gradient descent (SGD) is a cornerstone of stochastic optimization
and a workhorse of modern machine learning
\citep{bottou2018optimization}. At each iteration, it replaces the exact
gradient with a fresh stochastic estimate. Its simplicity and scalability
have made it a standard method across machine learning. For smooth
nonconvex objectives, the classical analysis gives an $O(K^{-1/2})$ guarantee
for the average squared gradient norm under unbiasedness and bounded
variance, with a suitably chosen stepsize
\citep{ghadimi2013stochastic}.

The bounded-variance assumption controls the mean-square magnitude of the
error in an individual stochastic gradient, uniformly over the points under
consideration. This uniform bound must cover the noisiest point in the region
and can therefore be much larger than the noise typically encountered along
the optimization trajectory. Classical SGD guarantees are consequently
sensitive to the magnitude of pointwise gradient noise.

But magnitude is not the only relevant feature of noise. In many problems,
the same data sample produces correlated stochastic gradients at nearby
points. Their difference can be accurate even when each gradient is noisy,
because noise shared across the two evaluations cancels. Plain SGD does not
explicitly exploit this coupling: each update uses a fresh estimate of the
absolute gradient. This suggests a complementary strategy---estimate how the
gradient changes, and use that change to carry an existing estimate along the
optimization path.

Transported SGD (TSGD) isolates this idea in its simplest form. Starting from
an initial gradient estimate, it estimates the gradient change between
consecutive iterates and adds that change to its current estimate. SGD
therefore estimates gradients afresh, whereas TSGD transports an initial
estimate by estimating gradient changes. Their error mechanisms differ: SGD
is sensitive to the magnitude of gradient noise, while TSGD is sensitive to
the variation of that noise along the trajectory. For example, with additive
noise shared across the two evaluations, the noise cancels exactly from the
gradient difference. Thus TSGD can tolerate large pointwise gradient
variance when the noise varies little along the trajectory and the initial
estimate is sufficiently accurate.

This distinction is captured by a similarity parameter $\delta$ that
controls the mean-square error in stochastic gradient changes. A small
$\delta$ is the key condition for accurate transport; it can arise when
component gradients are similar, but does not require a finite-sum objective.
TSGD's transport error therefore depends on noise variation rather than on a
uniform bound for pointwise noise magnitude. This does not make TSGD
independent of all noise: an inaccurate initial estimate can persist, and
transport errors can accumulate along the trajectory. These two error
sources shape its guarantees and distinguish its behavior from SGD.

The transport recursion itself is not new. It appears in SARAH and its
inexact variants, including SARAH-IN and iSARAH-IN
\citep{nguyen2017sarah,nguyen2021inexact}, and in the path-integrated
estimator underlying SPIDER \citep{fang2018spider}. It is also recovered by
PAGE when its refresh probability is set to $p=0$, and is related to PAGE-NS
\citep{li2021page,sokolov2022nonconvex}; the CGM estimator reduces to the same
recursion when its momentum parameter is $\beta=0$
\citep{jiang2026nonconvex}. Despite these appearances, transport is typically
embedded within a variance-reduction or momentum scheme rather than isolated
as a basic method. We argue that estimating gradient changes merits study as
a fundamental mechanism in its own right---a sister to fresh-gradient SGD
and momentum SGD (MSGD). We study this known mechanism on its own, not as a
claim to a new algorithm, to clarify what transport guarantees and when it
is effective.

Our contributions are fourfold:
\begin{itemize}
  \item We give a unified convergence analysis for TSGD that accounts for
  both inexact initialization and stochastic gradient changes.
  \item We characterize the separate effects of initialization and transport
  errors. A scalar quadratic model yields exact dynamics in all four cases
  formed by exact or inexact initialization and exact or stochastic
  transport.
  \item We compare SGD and TSGD under their respective noise assumptions and
  convergence guarantees, and use additive- and multiplicative-noise
  examples to identify regimes favoring each method.
  \item We connect the transport recursion to existing variance-reduced and
  momentum-based methods, separating transport from the refresh or momentum
  mechanisms combined with it.
\end{itemize}
